\documentclass[a4paper,11pt]{article}
\usepackage{fontspec,xeCJK}
\setmainfont{Times New Roman}
\setsansfont{Arial}
\setmonofont{Consolas}[Scale=0.87]
\setCJKmainfont{Yu Mincho}[BoldFont=Yu Mincho Demibold]
\setCJKsansfont{Yu Gothic}
\setCJKmonofont{Yu Gothic}
\usepackage{amsmath,amssymb,graphicx,booktabs,longtable,tabularx,array}
\usepackage[a4paper,top=23mm,bottom=23mm,left=23mm,right=23mm]{geometry}
\usepackage{xcolor,caption,fancyhdr}
\usepackage[unicode,colorlinks=true,linkcolor=black,citecolor=black,urlcolor=blue!55!black,bookmarksnumbered]{hyperref}
\hypersetup{pdftitle={OESCRの理論・実装と数値検証},pdfsubject={低圧プラズマの縮約衝突放射モデルと分光逆解析}}
\renewcommand{\contentsname}{目次}
\renewcommand{\figurename}{図}
\renewcommand{\tablename}{表}
\renewcommand{\refname}{参考文献}
\renewcommand{\abstractname}{要旨}
\numberwithin{equation}{section}
\setlength{\parskip}{3pt}
\setlength{\parindent}{1em}
\setlength{\tabcolsep}{5pt}
\renewcommand{\arraystretch}{1.25}
\linespread{1.13}
\emergencystretch=1em
\makeatletter
\setlength{\@fptop}{0pt}
\makeatother
\captionsetup{font=small,labelfont=bf,labelsep=quad}
\pagestyle{fancy}
\fancyhf{}
\fancyhead[L]{\small OESCR：理論・実装と数値検証}
\fancyhead[R]{\small 2026-10-27 実装記述}
\fancyfoot[C]{\thepage}
\setlength{\headheight}{15pt}
\newcolumntype{Y}{>{\raggedright\arraybackslash}X}
\newcommand{\code}[1]{\texttt{#1}}
\newcommand{\unit}[1]{\,\mathrm{#1}}
\newcommand{\nelec}{n_{\mathrm e}}
\newcommand{\Te}{T_{\mathrm e}}
\newcommand{\fe}{f_E}
\newcommand{\Aeff}{A^{\mathrm{eff}}}
\newcommand{\figfile}[1]{../figures/#1}
\title{\LARGE 低圧プラズマ発光分光のための\\縮約衝突放射・観測モデル OESCR\\[6pt]\large 理論的背景、実装方程式、逆解析および再現ベンチマーク}
\author{}
\date{2026年10月27日実装に基づく改訂版}
\begin{document}
\maketitle
\begin{abstract}
OESCRは、低圧プラズマの電子状態と粒子密度から発光スペクトルを計算し、複数の観測に共通する状態量を推定する研究用Pythonコードである。電子エネルギー分布に基づく衝突速度係数、選択した励起準位の定常収支、視線積分および分光器応答を一つの計算経路として結ぶ。利用者は、対象ガスの状態・反応・放射分岐を組み合わせ、温度、電子密度、ラジカル密度等の推定に使う観測と既知量を指定する。本稿は、この計算を支える方程式、近似の物理的意味、機能群の分担、数値解法と実装上の工夫を述べる。解析解との比較および生成スペクトルの回復試験では、積分光量の保存、比と絶対尺度の使い分け、共有状態の推定、多視線・広帯域観測への接続を確認した。Ar/O$_2$では温度・密度を良好に回復し、Ar/Cl$_2$では温度と密度の感度分離が診断精度を左右することを示した。現行の計算機能を分光モデルの構築、反応機構の比較、診断線と観測配置の設計に位置づけ、校正済み実測と独立温度診断による定量評価を次の検証段階として整理する。
\end{abstract}
\noindent\textbf{キーワード：} 発光分光、低温プラズマ、衝突放射モデル、電子エネルギー分布、逆問題、識別可能性
\clearpage
\section*{全体像：何を入力し、何を計算し、何を推定するか}
\addcontentsline{toc}{section}{全体像：機能群と計算フロー}
OESCRは、発光分光（optical emission spectroscopy; OES）と衝突放射（collisional-radiative; CR）計算を接続する、縮約された順・逆解析基盤である。電子エネルギー分布関数（electron energy distribution function; EEDF）と粒子密度を励起計算の入力とする。本稿のOESCRは当該コードを指し、単一の普遍的なプラズマ方程式の名称ではない。まず既知または仮定したプラズマ状態から測定値を予測し、その予測と観測のずれから、宣言した少数の状態量を推定する。

\begin{figure}[htbp]
\centering
\newcommand{\flowbox}[2]{\fbox{\parbox{0.90\linewidth}{\textbf{#1}\quad #2}}}
\flowbox{入力・モデルの定義}{圧力・ガス温度・組成、電子と供給粒子の密度、EEDF、状態・反応・放射データ、視線と分光器、測定CSV}
\\[2pt]$\downarrow$\\[2pt]
\flowbox{1. 読込・検査・静的準備}{YAMLとspecies packを展開。参照・単位・密度の所有者を検査し、格子・経路長・固有線形状を準備する。}
\\[2pt]$\downarrow$\\[2pt]
\flowbox{2. 励起動力学}{領域ごとのEEDF $\to$ 反応速度係数 $\to$ 線形定常CR $M\boldsymbol n^*=\boldsymbol b$ $\to$ 原子線・分子帯の固有発光。}
\\[2pt]$\downarrow$\\[2pt]
\flowbox{3. 観測演算}{視線積分 $\to$ 入射波長での透過率・単位換算 $\to$ 波長写像・LSF $\to$ gain・baseline $\to$ 画素幅で平均。}
\\[2pt]$\downarrow$\\[2pt]
\flowbox{4. 推定・出力}{観測との共分散で白色化した残差を最小化。名前付き物理量、予測・残差を保存。必要時に品質・SVD・profile・局所曲率を監査する。}
\caption{計算の全体フロー。逆解析では、推定パラメータを更新して段階2--4を反復する。ガス組成の生成、電力収支、電場、輸送を求める自己無撞着ループはこのフローに含まれない。}
\label{fig:overview}
\end{figure}

\begin{table}[htbp]
\centering\small
\begin{tabularx}{\linewidth}{p{27mm}YY}
\toprule
機能群 & 役割・主な出力 & 実装の所在 \\
\midrule
設定とデータ & 入力の正規化、種別データの合成、出典・hash & \code{io}, \code{data} \\
物理計算 & EEDF、速度係数、準位密度、発光、生成・消滅収支 & \code{physics}, \code{forward.emissivity} \\
幾何と装置 & 殻領域から視線への投影、装置単位でのスペクトル & \code{geometry}, \code{instrument} \\
順解析 & 固定構造のcompileと繰返し評価 & \code{forward.compiled}, \code{forward.model} \\
逆解析 & 測定誤差、目的関数、境界付き推定、任意の監査 & \code{inverse} \\
検証・報告 & 解析解との比較、回復誤差、保留窓の評価、根拠区分 & \code{analysis}, \code{scripts} \\
\bottomrule
\end{tabularx}
\caption{機能の分担。モジュール名はすべて\code{oescr}パッケージ内の相対名である。}
\end{table}

\noindent\textbf{読み方。} 第1--2節で低圧プラズマと励起方程式を導入し、第3節で測定値への変換、第4節で逆問題を説明する。第5節は物理項の変更と複数ガスへの拡張、第6節は再現ベンチマーク、第7節は利用範囲を扱う。数式中のエネルギーは原則eV、波長はnm、密度はm$^{-3}$とする。数式は物理量として記し、実際の設定キーとの対応は付録にまとめる。

\clearpage
{\linespread{1}\fontsize{9.5}{10.2}\selectfont\setlength{\parskip}{0pt}\tableofcontents}
\clearpage
\section{背景とモデルの位置づけ}
\subsection{低圧プラズマの発光診断で求めるもの}
半導体のエッチング、成膜、チャンバー洗浄に用いる容量結合プラズマ（CCP）と誘導結合プラズマ（ICP）では、電子が電場からエネルギーを受け取り、原子・分子の励起、解離、電離を生じる。電子の特徴的なエネルギーはeV程度に達し、重粒子はそれより低い温度に保たれることが多い。この非平衡状態では、励起準位への衝突供給と放射・消光による損失を計算することが、発光を解釈する基本となる\cite{fujimoto,lieberman}。

OESは光を受動的に観測し、多数の原子線と分子帯を同時に取得する方法である。例えばF I線の強度には、F原子密度、高エネルギー電子による励起、親分子からの解離励起、衝突消光、視線内分布、波長応答が寄与する。分光診断で粒子密度や電子状態を求めるには、これらの寄与を反応と光学応答に分けて評価する\cite{coburn,fuller}。

OESCRはこの評価を、選択した発光種・準位と実験の観測条件に合わせて構成する。対象とする時間区間には、準位の定常近似を適用できる区間を選ぶ。幾何には、軸対称の殻モデルや外部で定めた投影行列を用いる。採用する励起・消光・捕獲の項は、放射寿命、衝突頻度、光学的厚さをもとに選定する。この手順によって、数Paの放電から、NF$_3$洗浄の例で参照する約133 Paの条件まで、それぞれの実験に即したモデルを組み立てられる\cite{an2023}。

\subsection{衝突放射モデルの考え方}
CRモデルは、励起準位への衝突供給、準位間移行、放射遷移、衝突失活を粒子収支で表す。放射損失が支配的な低密度極限では、基底状態からの直接励起と放射が釣り合うコロナモデルに近づく。密度が上がるにつれて、準位間衝突や準安定状態からの段階励起が重要になり、必要な状態と反応経路を収支へ加える\cite{fujimoto}。

OESCRでは、電子密度、EEDF、衝突相手の密度を与えたときの準位収支を線形方程式として解く。逆解析では、その外側で温度や密度を更新して速度係数と収支を再評価する。この構成により、一回の準位計算の見通しを保ちながら、非線形な発光応答に基づく状態推定を行える。

\begin{table}[htbp]
\centering\small
\begin{tabularx}{\linewidth}{p{26mm}YYY}
\toprule
モデル & 主に解く量 & 主な用途 & OESCRとの接点 \\
\midrule
LTE・Boltzmann分布 & 温度で定まる準位占有率 & 熱平衡に近い発光の解析 & 非平衡収支と比較する基準 \\
コロナモデル & 直接励起と放射の釣合い & 放射損失が支配する希薄系の線強度 & CRの採用項を限定した極限 \\
詳細CRモデル & 多数準位・多価数の占有率 & 機構を詳しく含めた分光診断 & 対象準位を選ぶ縮約の出発点 \\
0次元化学モデル & 化学種密度、粒子・電力収支 & 組成、解離率、運転条件の評価 & 供給粒子密度や電子状態の入力源 \\
流体・輸送モデル & 密度、流束、場の空間分布 & 放電構造と装置条件の評価 & 密度・温度分布を殻状態へ、観測配置を投影重みへ反映 \\
PIC/MCC & 荷電粒子の運動、衝突、電場 & 非局所加熱、シース、電子動力学 & 電子分布等を発光計算へ入力 \\
Boltzmann（BOLSIG+） & 二項近似の電子速度分布 & 指定した均一場・$E/N$下の輸送・反応係数 & 表形式EEDFと率積分の比較 \\
OESCR & 準位密度、測定スペクトル、指定した推定量 & 分光解釈、診断設計、機構・校正の感度比較 & 物理状態と観測値を結ぶ計算 \\
\bottomrule
\end{tabularx}
\caption{各モデルが担当する問題と、OESCRへの入出力の接点\cite{fujimoto,lieberman,hagelaar,arellano}。外部モデルの結果は、対応する密度・EEDF・幾何をcaseへ与えて利用する。}
\label{tab:models}
\end{table}

\subsection{利用目的に応じた計算の組み立て}
発光機構を調べる場合は、基底励起、準安定からの励起、消光等を選んで順解析し、準位収支と線強度の変化を比較する。測定から状態量を求める場合は、既知の供給密度と校正条件を固定し、観測が支える温度・密度を推定対象にする。診断設計では、異なる励起感度を持つ線、参照種、視線を加えたときのJacobianやprofileの変化を調べる。

複数ガスへの展開も同じ考え方に従う。各ガスの状態・反応データを用意し、共通の電子状態、混合相手による消光、観測応答を設定することで、同一の解析手順を使える。対象条件の選定とデータの根拠を明示することが、計算結果を実験の解釈へつなぐ要点となる。

\section{プラズマ状態、電子衝突および準位収支}
\subsection{領域ごとの入力状態と計算対象}
領域$z$の状態を、電子密度$\nelec(z)$、エネルギー空間の電子確率密度$\fe(E;z)$、外部供給密度$N_s(z)$、解く励起準位密度$n_i^*(z)$で表す。ガス圧力$p$とガス温度$T_g$から親ガスの密度を
\begin{equation}
N_{\mathrm{gas}}=\frac{p}{k_{\mathrm B}T_g},\qquad N_s=x_sN_{\mathrm{gas}}
\label{eq:gas}
\end{equation}
とする。$x_s$は入力のモル分率である。親ガスは圧力・温度・分率から、ラジカル、準安定状態、残留ガスは領域別の密度指定から構成する。例えばCF$_4$/O$_2$/Arでは、親ガス組成に加えてF、CF、CF$_2$、Oの密度を与える。この入力密度を測定値または外部の化学モデルから定めることで、対象とする励起動力学に計算を集中できる。

各密度は一つの入力欄が所有する。電子、親ガス、ラジカル、準安定状態と、CRで解く準位の対応を検査し、同一種の二重指定は入力エラーとする。\code{solve: false}の準位は既知の密度として扱い、その準位からの発光と、解く下準位への放射供給を計算できる。これにより、独立に測定した準安定・励起準位密度を発光計算へ組み込める。

\subsection{電子エネルギー分布関数と温度の意味}
本コードは、電子エネルギー$E$を変数とする確率密度$\fe(E)$を用いる。$\fe(E)\,dE$は、電子が$E$から$E+dE$の区間に存在する割合である。速度空間分布や電子エネルギー確率関数（EEPF）の資料を使う場合は、この定義へ換算した分布を入力する。
\begin{equation}
\int_0^\infty \fe(E)\,dE=1,\qquad [\fe]=\mathrm{eV}^{-1},\qquad
\langle E\rangle=\int_0^\infty E\fe(E)\,dE .
\label{eq:pdf}
\end{equation}
Maxwell分布の入力$\Te$は、$k_{\mathrm B}T$をeVで表した量であり、十分な積分範囲では
\begin{equation}
f_{\mathrm M}(E;\Te)=\frac{2}{\sqrt\pi}\frac{\sqrt E}{\Te^{3/2}}
\exp\!\left(-\frac{E}{\Te}\right),\qquad \langle E\rangle=\frac32\Te .
\label{eq:maxwell}
\end{equation}
励起閾値より上の分布が速度係数を左右するため、平均エネルギーが同じでも、高エネルギー成分が異なれば発光は異なり得る\cite{hagelaar}。

Druyvesteyn型は、実装の\code{te\_eV}を尺度$E_D$と読み替えて
\begin{equation}
f_{\mathrm D}(E;E_D)=\frac{2\sqrt E}{E_D^{3/2}\Gamma(3/4)}
\exp\!\left[-\left(\frac E{E_D}\right)^2\right],\qquad
\langle E\rangle=E_D\frac{\Gamma(5/4)}{\Gamma(3/4)}
\label{eq:druy}
\end{equation}
と書ける。式\eqref{eq:druy}の係数と平均値は無限区間の解析値であり、コードは有限格子で数値的に正規化する。$E_D$は分布の尺度である。Maxwell分布との比較には、両分布の平均エネルギーを式\eqref{eq:maxwell}、\eqref{eq:druy}から求める。

二温度Maxwell分布は
\begin{equation}
f_{\mathrm{bi}}(E)=(1-\alpha)f_{\mathrm M}(E;T_c)+\alpha f_{\mathrm M}(E;T_h),
\qquad 0\leq\alpha\leq1,\quad T_h\geq T_c>0
\label{eq:bimax}
\end{equation}
である。$\alpha$は高温成分に属する電子の粒子割合で、$T_c$、$T_h$とともに指定する。二温度表現により、少数の高エネルギー電子が励起率へ及ぼす影響を調べられる。外部で得た分布を直接用いる場合は表形式EEDFを選び、内部格子への線形補間、表の範囲外のゼロ補完、再正規化を経て、速度係数と発光を順解析する。

数値計算では、実際に使う格子$[E_{\min},E_{\max}]$上で各分布を再正規化する。積分範囲は励起に寄与する高エネルギー側まで含める。上端付近の確率質量と分布振幅を診断し、必要に応じて$E_{\max}$を拡大して速度係数の収束を確認することで、裾の切断が発光へ及ぼす影響を評価する。

\subsection{電子衝突速度係数と断面積}
反応$r$の断面積$\sigma_r(E)$とEEDFから、速度係数を
\begin{equation}
k_r[\fe]=\int_0^\infty \sigma_r(E)\sqrt{\frac{2eE}{m_{\mathrm e}}}\,\fe(E)\,dE,
\qquad [k_r]=\mathrm{m^3\,s^{-1}}
\label{eq:rate}
\end{equation}
で得る。$e$はeVをJへ換算する素電荷、$m_{\mathrm e}$は電子質量である。エネルギーの単位換算は速度$\sqrt{2eE/m_{\mathrm e}}$に含まれ、$\fe(E)\,dE$は無次元の確率を与える。断面積を共有エネルギー格子へ線形補間し、台形則で積分する。一定の速度係数や、閾値を持つ解析的な代用断面積も選べるため、解析解との比較や機構の感度調査を簡潔に構成できる。実ガスの定量計算では、対象過程に対応する断面積と出典を選定する。

断面積CSVは\code{energy\_eV,sigma\_m2}の二列とし、エネルギーの単調増加、有限性、断面積の非負性を検査する。積分処理では表の外をゼロで補完し、その範囲のEEDF確率を未被覆量として記録する。利用する電子分布の裾に対応したエネルギー範囲の表を用意することが、率計算の入力条件となる。物理閾値から始まる表では、先頭行の断面積をゼロとして\code{zero\_below\_threshold}を宣言すると、閾値未満を既知のゼロ領域として被覆に含める。

断面積の出典、過程名、エネルギー範囲、SHA-256を記録するため、LXCat等から抽出した入力を追跡できる\cite{lxcat}。被覆率は電子確率分布と表のエネルギー範囲の対応を示す。速度係数の精度を調べるときは、被覆率に加え、積分核$\sigma v\fe$、励起閾値、および積分範囲の拡大に対する率の変化を確認する。

\subsection{定常線形CR方程式の組立て}
解く準位$i$への外部供給を$S_i$、追跡する準位$j$から$i$への一次移行頻度を$\nu_{j\to i}$、追跡外への損失を$\nu_{i\to\mathrm{out}}$とする。
\begin{equation}
\frac{dn_i^*}{dt}=S_i+\sum_{j\ne i}\nu_{j\to i}n_j^*
-\left(\sum_{j\ne i}\nu_{i\to j}+\nu_{i\to\mathrm{out}}\right)n_i^* .
\label{eq:balance}
\end{equation}
電子分布と衝突相手密度を固定すれば、すべての$\nu$は既知のs$^{-1}$である。定常化により
\begin{equation}
M\boldsymbol n^*=\boldsymbol b,\qquad
M_{ii}=\sum_{j\ne i}\nu_{i\to j}+\nu_{i\to\mathrm{out}},\qquad
M_{ij}=-\nu_{j\to i}\ (i\ne j),\qquad b_i=S_i
\label{eq:matrix}
\end{equation}
を解く。外部密度$N_s$からの供給は$S_i\mathrel{+}=\nu_{s\to i}N_s$として右辺へ入る。未知準位を出発点とする過程は、対角の損失と、追跡する到達準位への非対角供給を同時に加える。各過程の行列寄与と生成・損失収支を同じ記録から構築するため、特定の項が結果を支配する理由を追跡できる。

\begin{table}[htbp]
\centering\small
\begin{tabularx}{\linewidth}{p{31mm}p{38mm}p{24mm}Y}
\toprule
過程 & 有効移行頻度$\nu$ & 係数の単位 & 代表的用途 \\
\midrule
電子衝突 & $\nelec k_e[\fe]$ & m$^3$ s$^{-1}$ & 基底・準安定からの励起、宣言した準位移行 \\
一次過程 & $k_1$ & s$^{-1}$ & 非放射損失、指定した一次移行 \\
二体過程 & $k_2N_M$ & m$^3$ s$^{-1}$ & 中性粒子による消光 \\
三体過程 & $k_3N_MN_N$ & m$^6$ s$^{-1}$ & 外部衝突相手を伴う有効移行・損失 \\
放射過程 & $\Aeff_{u\ell}$ & s$^{-1}$ & 放射損失、追跡下準位へのcascade \\
壁損失 & $\nu_{w,i}$ & s$^{-1}$ & 指定した準位の有効境界損失 \\
\bottomrule
\end{tabularx}
\caption{実装済みの移行・損失項。二体・三体過程では、衝突相手の外部密度を既知係数として線形性を保つ。未知準位を衝突相手に指定する入力は反応式の検査でエラーとする。}
\label{tab:terms}
\end{table}

\subsection{二準位の例：式とコードの対応}
外部基底粒子から準位1、2へ$S_1,S_2$が供給され、準位2から1へ放射cascade $A_{21}$があるとする。$D_i$を各準位の全損失頻度とすると
\begin{equation}
\begin{pmatrix}D_1&-A_{21}\\0&D_2\end{pmatrix}
\begin{pmatrix}n_1^*\\n_2^*\end{pmatrix}
=\begin{pmatrix}S_1\\S_2\end{pmatrix},\qquad
n_2^*=\frac{S_2}{D_2},\quad
n_1^*=\frac{S_1+A_{21}S_2/D_2}{D_1}.
\label{eq:twolevel}
\end{equation}
$D_2$には$A_{21}$自体も含む。消光を加えれば該当する$D_i$が増え、cascadeの除外は上準位の損失と下準位の供給の双方を変える。単一準位の場合は
\begin{equation}
n_u^*=\frac{\nelec N_s k_{s\to u}}
{\sum_\ell\Aeff_{u\ell}+\sum_M k_{q,uM}N_M+\nu_{w,u}+\nu_{\mathrm{other},u}}
\label{eq:single}
\end{equation}
となり、密度と温度以外の入力が発光へどう入るかを明瞭に読める。

この準定常解は、励起系の緩和が入力状態の変化に十分速く追従する条件で用いる。緩和時間は単独準位で概ね$D_u^{-1}$、結合系で行列の緩和固有モードを目安に評価する。利用者はこれを電子状態や供給密度の変動時間と比較し、解析する時間区間を選ぶ。時間平均スペクトルへ適用する場合は、発光の入力依存性に由来する平均化の影響も評価する。時間発展と輸送の連成は第\ref{sec:outlook}節の拡張課題で扱う。

\subsection{放射捕獲、分岐、壁損失}
放射捕獲は、指定した光学的厚さ$\tau$に対する簡略slab脱出係数で表す。
\begin{equation}
\beta(\tau)=\frac{1-\exp(-\tau)}{\tau},\qquad \beta(0)=1,\qquad
\Aeff_{u\ell}=\beta_{u\ell}A_{u\ell}.
\label{eq:escape}
\end{equation}
この実装は、指定した光学的厚さ$\tau$または脱出係数$0\leq\beta\leq1$から、有効放射率$\Aeff$を求めるslab近似である。同じ$\Aeff$を準位収支と発光の双方へ用いることで、捕獲による寿命の変化と外へ逃げる光子の生成を一貫して扱う。数値的には$\tau\leq10^{-12}$で1とし、slabの$\beta$を$[10^{-6},1]$に制限する。例えば放射損失だけを持つ単一経路では、$\beta$の減少が占有率の増加を伴い、定常状態の全脱出光子数が保たれる極限が得られる。光学的厚さの値は、対象線の吸収と観測条件に基づいて与える。

原子データは波長と放射率の集合として管理し、各ケースの衝突過程と組み合わせる。例えばArの2p1は667.7、750.4 nm、2p6は763.5、800.6、922.4 nmへの分岐を持つ。観測範囲外も含む分岐集合を$D_u$へ入れることで、上準位寿命と観測線の光子収率を求める。NISTの波長と$A$値を固定した種別packは、この集合を複数のケースで再利用するために有用である\cite{nist}。

壁損失は
\begin{equation}
\nu_{w,i}=\gamma_i f_{\mathrm{flux}}\frac{\bar v_i}{L},\qquad
\bar v_i=\sqrt{\frac{8k_{\mathrm B}T_g}{\pi m_i}},\qquad
L=\frac VA=\frac{RH}{2(H+R)}
\label{eq:wall}
\end{equation}
とする。$R,H$は円柱容器の半径・高さで、$L$の直接指定も可能である。$\gamma_i$は状態別または種別の損失確率、$m_i$は質量である。この項は、指定状態に対する境界損失を一次の消滅頻度として準位収支へ加える。\code{thermal\_flux\_factor}の既定値は互換性のため1であり、等方的な平均速さと通常の熱流束$N\bar v/4$を使う定義では$1/4$を明示する。非ゼロの壁損失を用いる状態には質量を指定し、$\gamma_i$と$L$を対象の壁面条件に合わせて定める。

\subsection{準位解と格子収束の確認}
準位系には密行列の直接解法を用いる。各未知準位を放射・衝突・壁損失で釣り合わせ、一意な定常解を持つ反応集合を与える。特異系は通常の予測時にもエラーとし、利用者が損失経路や固定密度を点検できるようにする。分子帯だけを計算するケースでは、空のCR系を使用できる。監査では行列のrank・条件数、相対残差$\|M\boldsymbol n^*-\boldsymbol b\|_2/\|\boldsymbol b\|_2$、および負の準位密度量を記録する。実装は負値を大きさによらずゼロへ置き換えるため、置換前の解と残差を用いて計算の受理を判断する。

格子収束の確認では、エネルギー格子と波長格子を個別に細分化し、同じ検出画素上で予測の相対L2差を比較する。この手順により、速度係数積分と線形状・画素積分に必要な分解能を決められる。反応集合と断面積の妥当性は、生成・損失収支の比較および第\ref{sec:benchmarks}--\ref{sec:outlook}節の検証手順で評価する。

\section{固有発光から測定スペクトルへの変換}
\subsection{原子線の発光係数}
上準位$u$から下準位$\ell$への放射について、等方発光を仮定した体積・立体角・波長当たりの発光係数を
\begin{equation}
j_{u\ell,\lambda}=a_{u\ell}\frac{hc}{4\pi\lambda_{u\ell}}
n_u^*\Aeff_{u\ell}\,\phi_{u\ell}(\lambda),\qquad
\int\phi_{u\ell}(\lambda)\,d\lambda=1
\label{eq:line}
\end{equation}
とする。$a_{u\ell}$は明示的な\code{amplitude\_scale}で既定値1、$[j_\lambda]=\mathrm{W\,m^{-3}\,sr^{-1}\,nm^{-1}}$である。光子エネルギーに使う波長はmへ変換する。$1/(4\pi)$があるため、後段の放射輝度を全立体角で積分した発光パワーと混同しない。

固有原子線形は、指定した中心とFWHMに対するGaussianである。
\begin{equation}
\phi(\lambda)=\frac{1}{\sqrt{2\pi}s_\lambda}
\exp\!\left[-\frac{(\lambda-\lambda_0)^2}{2s_\lambda^2}\right],\qquad
s_\lambda=\frac{\mathrm{FWHM}}{2\sqrt{2\ln2}}.
\label{eq:gaussian}
\end{equation}
Doppler幅、Stark幅、圧力幅を温度や密度から自動的に計算する機能ではない。線形は観測区間と独立に正規化する。観測窓で線の片側が欠けたとき、その失われた光を残った区間へ再配分しない。

同一上準位の光学的に薄い二分岐を完全積分した強度比は、応答補正後には
\begin{equation}
\frac{J_{u\ell_1}}{J_{u\ell_2}}
=\frac{a_{u\ell_1}A_{u\ell_1}/\lambda_{u\ell_1}}
{a_{u\ell_2}A_{u\ell_2}/\lambda_{u\ell_2}}
\label{eq:branchratio}
\end{equation}
となり、占有率$n_u^*$は消える。分岐線は波長応答や放射分岐の整合性の確認には有用だが、線の本数だけ独立なEEDF情報が増えるわけではない。本稿でいう励起kernelは、断面積から上準位占有率を経て発光へ至る応答経路を区別する単位であり、数値的な感度の独立性まで保証する語ではない。

\subsection{物理的な分子帯と経験的な分子帯}
\code{electron\_impact\_photon\_band}は、供給粒子密度$N_s$、速度係数$k_b$、反応当たり光子収率$Y_b$、分岐割合$B_b$を使って
\begin{equation}
j_{b,\lambda}=\frac{hc}{4\pi\lambda}\nelec N_s k_b[\fe]Y_bB_b\psi_b(\lambda),
\qquad \int\psi_b(\lambda)\,d\lambda=1
\label{eq:band}
\end{equation}
を与える。帯域内の各波長で光子エネルギーを換算するため、原子線と同じ放射パワー基準で加算できる。$Y_b$は明示入力であり、コードが励起状態の消光や振動分布から導出するものではない。帯の形はGaussianまたは表形式profileで指定し、表形式は元の全波長支持域で正規化する。

一方、\code{effective\_excitation\_band}は$a\nelec N_sk_b\psi_b$、\code{effective\_density\_band}は$cN_s\psi_b$に比例する経験的出力である。これらは事象率または経験係数で尺度づけた形で、放射輝度と同じ次元とは限らない。使用には\code{emission\_mode: empirical}を明示し、絶対強度推定には使わない。現行のNF$_3$/Ar、Cl$_2$/Ar例にはこの経験的成分がある。分子帯の回転・振動準位を連立CRで解く機能や、帯形から回転温度を自動推定する機能は含まれない。

\subsection{視線積分と軸対称の殻モデル}
半径方向に等幅の殻領域を置き、殻$z$の発光を一定と近似する。同一高さの視線$m$のimpact parameterを$b_m$、殻の内外半径を$r_z,r_{z+1}$とすると
\begin{equation}
I_{m,\lambda}=\sum_zW_{mz}j_{z,\lambda},\qquad
W_{mz}=2\left[\sqrt{(r_{z+1}^2-b_m^2)_+}-\sqrt{(r_z^2-b_m^2)_+}\right],
\label{eq:chord}
\end{equation}
ただし$(x)_+=\max(x,0)$である。$W$の単位はm、$I_\lambda$は放射輝度$\mathrm{W\,m^{-2}\,sr^{-1}\,nm^{-1}}$となる。各殻のCR計算は独立で、殻間拡散や輸送による結合はない。多視線の逆解析は、この投影を含む少数殻の状態推定であって、任意の二・三次元場の再構成ではない。

外部行列を使う\code{user\_field}で別の幾何を接続できる。\code{chordavg}は全1の行列で領域を加算する簡略機能であり、それ自体は実長を持つ視線積分ではない。\code{asym\_lowrank}は視線ごとに$1+a_m\cos\theta_m$を掛ける補正で、角度方向のトモグラフィではない。絶対強度に用いる場合は、選択した投影重みの物理的長さと観測体積の意味を明示する。

\subsection{装置応答の順序と単位}
入射放射輝度$I$に対し、透過率$T(\lambda)$、次元換算$\mathcal C$、波長写像$\mathcal S_\delta$、線広がり関数（line spread function; LSF）$K$、固定gain $g$、baseline $b$を順に適用する。画素$i$の予測は
\begin{equation}
F_i=\frac{1}{\Delta\lambda_i}\int_{\lambda_{i-1/2}}^{\lambda_{i+1/2}}
\left\{g\left[K*\mathcal S_\delta\!\left(\mathcal C[TI]\right)\right](\lambda)
+b(\lambda)\right\}\,d\lambda .
\label{eq:instrument}
\end{equation}
光子数換算と透過率は入射した物理波長で計算する。波長依存の透過率とLSFの畳み込みは一般に交換できず、波長写像後の見かけの波長から光子エネルギーを再計算しない。現在の波長shiftの実装は$\mathcal S_\delta q(\lambda)=q(\lambda+\delta)$で、正の\code{wavelength\_shift\_nm}はピークを短波長側へ移す。

装置LSFはGaussianまたはVoigtで、離散kernelの総和を1にして波長方向へFFT畳み込みを行う。現行のkernelは波長に対して並進不変で、有限支持域のVoigt翼は切断後に再正規化する。透過率は定数または表、baselineは定数または多項式である。観測端での切断やLSF翼の影響は別途点検する。

\begin{table}[htbp]
\centering\small
\begin{tabularx}{\linewidth}{p{35mm}p{48mm}Y}
\toprule
出力基準 & 変換 & 出力単位 \\
\midrule
相対装置信号 & 絶対校正を主張しない応答 & arb \\
分光放射輝度 & $I_\lambda$を維持 & W m$^{-2}$ sr$^{-1}$ nm$^{-1}$ \\
収集分光パワー & $P_\lambda=I_\lambda A_c\Omega f_v$ & W nm$^{-1}$ \\
光電子スペクトル & $N_{e,\lambda}=P_\lambda t\eta/(hc/\lambda)$ & photoelectron nm$^{-1}$ \\
\bottomrule
\end{tabularx}
\caption{装置出力の基準。$A_c$は収集面積、$\Omega$は立体角、$f_v$はviewing factor、$t$は積分時間、$\eta$は量子効率。現行の量子効率はスカラー入力である。baselineは変換後の出力単位で加算する。}
\end{table}

検出画素には中心波長の点値ではなく、式\eqref{eq:instrument}の\textbf{画素平均}を出力する。したがって、光電子スペクトルを画素当たりの光電子数へ直すには画素幅を掛ける。細い原子線を内部の細格子で保持し、画素境界を含む積分を最後に行うため、粗い測定画素の間へ未分解線が落ちて消える問題を避けられる。この処理は低分解能分光器で線強度と帯形を比較する際に重要である。

\section{逆解析、測定誤差および識別可能性}
\subsection{推定する状態量と観測の共有}
推定対象$\boldsymbol\theta$は、温度、電子密度、選択した外部粒子密度など、設定ファイルの明示的な経路へ結びつける。正の密度等には$x=\log_{10}\theta$を使え、上下限もその座標へ変換する。報告値は名前付きの物理単位へ戻す。

複数の装置、視線、繰返し測定を同じプラズマ状態へ接続できる。観測IDと予測IDを分離するため、三回の観測を一つの予測スペクトルへ対応させることも可能である。現行の一回のfitは、一つの\code{condition\_id}と一つの\code{state\_group\_id}を扱う。異なる電力・圧力条件を同時に推定する階層モデルとは区別する。

\subsection{残差と共分散}
選択した画素の観測を$\boldsymbol y$、予測を$\boldsymbol F(\boldsymbol\theta)$とする。測定共分散$C=LL^\mathsf T$に対し、
\begin{equation}
\boldsymbol r_{\mathrm{data}}(\boldsymbol\theta)
=L^{-1}\left[\boldsymbol F(\boldsymbol\theta)-\boldsymbol y\right]
\label{eq:whiten}
\end{equation}
と白色化する。独立誤差なら$r_i=(F_i-y_i)/\sigma_i$に戻る。共分散は有限・対称・正定値を要し、明示的な逆行列ではなくCholesky因子で解く。測定$\sigma$も共分散もない相対fitでは、観測の標準偏差等による便宜的尺度を使う。この場合のcostを、校正済みの統計的$\chi^2$と同じ意味で扱うことはできない。

同じ校正を共有する観測には、共通の乗算誤差がある。校正group $g$に属する観測$i$の相対標準不確かさを$s_i$、観測値を使った固定ベクトルを$\boldsymbol u_g$として
\begin{equation}
C=C_{\mathrm{meas}}+\sum_g\boldsymbol u_g\boldsymbol u_g^\mathsf T,
\qquad (u_g)_i=\begin{cases}s_iy_i,&i\in g,\\0,&i\notin g\end{cases}
\label{eq:calcov}
\end{equation}
とする。予測値で共分散を毎回更新するのではなく、観測を尺度基準として固定する。共通誤差を独立画素誤差とみなしてしまうと、画素数や繰返し数に比例して精度が向上したように見えるため、この区別が重要になる。

測定の画素状態を保持し、通常残差には有効な画素を使う。飽和画素には下限値のメタデータを要求するが、現行の目的関数では有効画素以外を除外し、飽和の打切り尤度や一側制約を評価する機能は実装していない。入力adapterも、校正済みCSVの単位、誤差、状態、出典を検査して接続するもので、raw countsの校正や平滑化を暗黙に行わない。

\subsection{目的関数、波長選択、補助パラメータ}
推定は、測定由来の残差、事前分布、殻方向平滑化を組み合わせた
\begin{equation}
\widehat{\boldsymbol\theta}=\mathop{\arg\min}_{\boldsymbol\theta\in\Theta}
\frac12\left(\|\boldsymbol r_{\mathrm{data}}\|_2^2+
\|\boldsymbol r_{\mathrm{prior}}\|_2^2+
\|\boldsymbol r_{\mathrm{smooth}}\|_2^2\right)
\label{eq:objective}
\end{equation}
で行う。全スペクトル、選択窓内の形、baseline補正後の面積・peak・比、名前付きfeatureの共分散を使える。面積や比の残差は所定の尺度で無次元化してから白色化する。同じfeatureを共分散項とスカラー重み項に二重に入れない。ただし、同じ画素から作ったfeatureと全画素残差を別々に足した目的関数は、自動的に独立な尤度にはならない。採用した重みの意味と情報の重複を吟味する。

通常の\code{full\_spectrum}は全有効画素を使う。\code{window\_union}を選ぶと、物理的に対応づけたfit窓の和集合だけを使える。\code{score}窓は推定後の予測確認へ保留し、fit窓との区間重複を禁止する。弱い窓を閾値で除く機能はあるが、共分散で名前を指定したfeatureが欠落した場合には、残差の長さを黙って変えずエラーとする。

相対fitでは、gain、offset、必要なら波長方向のtiltを内側で推定できる。これらも最終残差と同じ有効画素と白色化を使う。スペクトルのみなら重み付き線形最小二乗、featureや事前項を含む場合は対応する内側評価を行う。固定装置gain・baselineは順解析で一度だけ適用する。外側の最適化は、必要に応じたdifferential evolutionの後、境界付き\code{least\_squares}で局所精密化する。

\subsection{四つの推定モードとactinometry}
\begin{table}[htbp]
\centering\small
\begin{tabularx}{\linewidth}{p{35mm}YY}
\toprule
モード & 主な情報 & 定量解釈に必要な条件 \\
\midrule
\code{relative\_shape} & 相対形状、制約された空間変化 & gainや供給密度との尺度縮退を確認する \\
\code{ratio\_diagnostic} & 有効な窓面積・peak比 & 波長間応答と比の動力学を与える \\
\code{actinometry} & 同じ比目的関数 & actinometerと対象の励起・消光・組成の前提を別途支持する \\
\code{calibrated\_absolute} & 絶対尺度の測定 & 物理発光、校正基準の一致、固定gain等、測定誤差、供給密度と損失の独立制約 \\
\bottomrule
\end{tabularx}
\caption{モードは利用目的と解釈の契約を明示する。選択した名前だけで、温度や密度の観測可能性が得られるわけではない。}
\end{table}

actinometryは、既知の参照粒子Aと対象Xの発光比から密度比を推定する考え方である\cite{coburn,fuller}。視線内で一様な状態、同じ電子密度・光路、光学的薄さ、直接励起、無視できる消光を仮定した単純なコロナ極限では
\begin{equation}
\frac{I_X}{I_A}=\frac{N_X}{N_A}\frac{k_X}{k_A}
\frac{B_X}{B_A}\frac{\lambda_A}{\lambda_X}\frac{T(\lambda_X)}{T(\lambda_A)}
\label{eq:actino}
\end{equation}
と書ける。$B$は放射分岐割合である。波長応答を補正した強度なら最後の透過率比は除く。実際には段階励起、解離励起、消光、EEDF変化の採否を確認しなければならない。OESCRの\code{actinometry}は、この前提を明示して同じ比の最適化を利用する。別の化学組成solverを起動するモードではない。

絶対モードでは自動gain・offset・tiltや装置尺度の同時推定を禁止し、測定CSVと装置の基準・単位・校正参照を照合する。しかし、$I\propto\nelec N_s k(\Te)$の$\nelec$と$N_s$をともに自由にすれば、校正済みでも積しか分からない場合がある。絶対校正の必要条件と、パラメータの識別条件は分けて判定する。

\subsection{SVD、profile、局所不確かさの読み方}
局所識別性は、事前分布や平滑化を除いた測定残差だけから、最適化座標$\boldsymbol x$に対する有限差分Jacobian
\begin{equation}
J_{ij}=\frac{\partial r_{\mathrm{data},i}}{\partial x_j},\qquad
J=U\,\operatorname{diag}(s_1,\ldots,s_p)V^\mathsf T
\label{eq:svd}
\end{equation}
を作って調べる。小さい特異値に対応する$V$の方向は、観測をほとんど変えずに動かせるパラメータの組合せである。条件数$s_{\max}/s_{\min}$とrankは座標尺度・局所点・閾値に依存する。full rankでも実用上の精度を保証しない。

明示したnuisanceパラメータのJacobianを$J_\eta$、targetを$J_\theta$とすると、
\begin{equation}
J_{\theta,\perp}=(I-Q_\eta Q_\eta^\mathsf T)J_\theta
\label{eq:project}
\end{equation}
で、nuisanceが説明できる測定方向を除く。$Q_\eta$はSVDによる列空間の正規直交基底である。内側で既にprofileしたgain等は二度数えない。事前分布を入れてrankを増やしたように見せることも避ける。

非線形性は、対象パラメータ$\psi$を固定して残り$\boldsymbol\lambda$を再推定する
\begin{equation}
D_p(\psi)=\min_{\boldsymbol\lambda}\|\boldsymbol r_{\mathrm{data}}(\psi,\boldsymbol\lambda)\|_2^2
\label{eq:profile}
\end{equation}
でも調べられる。局所的で境界付きのprofileであり、score窓、事前項、平滑化を除く。狭い閉じた谷か、広い谷か、境界へ開いているかを確認する。$\Delta D=1$や3.84は追加条件下の漸近的目安であり、現行の図から校正済み信頼区間と呼ぶことはしない。

Laplace近似は全目的関数の局所曲率から共分散を近似する。事前・平滑化・指定共分散に条件づけられ、断面積、欠落反応、モデル不一致の不確かさを周辺化した完全な事後分布ではない。これらの監査は\code{--audit}等で明示的に実行し、毎回の順解析反復には付けない。最適化の\code{success}、数値品質、実用的識別性、外部検証は、それぞれ別の判定である。

\section{物理項の改訂と複数ガスへの展開}
\subsection{物理項を変更する三つの方法}
OESCRでは、改訂の内容に応じて、データと設定の変更、関数形を置き換えるpluginの追加、連立方程式を拡張する計算核の開発を使い分ける。日常的なモデル改訂の中心は、状態・反応・遷移・分子帯を記したYAMLと、断面積CSVの更新である。

例えばO 777 nmの消光係数を改訂する作業では、対応する二体過程の係数を更新する。放射分岐の追加では遷移集合を更新し、励起断面積の変更では反応が参照するCSVを差し替える。供給と損失を同じ行列組立て規則で評価するため、これらの改訂は準位寿命と発光へ一貫して反映される。変更前後の準位収支とスペクトルを比較すれば、その物理項が観測へ及ぼす影響を調べられる。

\begin{table}[htbp]
\centering\small
\begin{tabularx}{\linewidth}{p{31mm}p{42mm}Y}
\toprule
改訂内容 & 実装済みの編集箇所 & 改訂時に確認する物理量 \\
\midrule
励起断面積 & 反応のCSVまたは\code{cross\_section\_file} & 状態・閾値・断面積単位・energy範囲・出典 \\
放射寿命・分岐 & species packの\code{transitions} & 全分岐による損失、air/vacuumの波長基準 \\
消光・一次失活 & \code{two\_body}/\code{first\_order}反応 & 衝突相手の入力密度と係数単位 \\
三体過程 & \code{three\_body}反応 & 二つの外部衝突相手の密度、係数の定義 \\
壁損失 & \code{wall}の$\gamma,L,f_{\rm flux}$ & 質量・熱流束定義・有効長 \\
放射捕獲 & 遷移の\code{trapping} & slab脱出係数、入力光学厚さ$\tau$ \\
分子帯 & \code{bands}の率・収率・profile & 光子収率、出力単位、profile支持域 \\
観測条件 & 視線幾何と\code{instruments} & 光路、透過率、LSF、画素、校正基準 \\
関数形 & 対応するplugin registry & 単位、正規化、解析解・保存則 \\
\bottomrule
\end{tabularx}
\caption{物理項ごとの改訂経路。既存の反応形式と観測演算を利用する変更は、該当するデータ・設定へ局所化できる。}
\label{tab:edits}
\end{table}

関数形を変更する場合は、EEDF、reaction rate、band emission/profile、trapping、wall loss、geometry、throughput、LSF、baselineのregistryへpluginを登録する。各pluginは選択キー\code{kind}、説明、入力schema、実行関数を備える。EEDFなら\code{build\_pdf}、反応率なら\code{rate}、壁損失なら\code{loss\_rate}を実装する。例えば、壁損失を別の有効長モデルへ変更し、順解析と逆解析から共通して利用することができる。

現在の計算核は、外部衝突相手密度を与えた線形定常CRと、領域ごとの独立計算を採用する。この構造の中ではデータ改訂とplugin追加でモデルを展開できる。時間発展、領域間輸送、未知密度同士の積を含む化学連成は、連立方程式と解法を拡張する開発に相当する。これらの位置づけと実測検証の順序を第\ref{sec:outlook}節にまとめる。

\subsection{species packによる複数種の組立て}
species packは、状態、反応、遷移、分子帯、出典をまとめる再利用単位である。例えばArの原子データとO 777 nmの縮約データは、次の入力で一つのcaseへ組み込める。
\begin{quote}\small\ttfamily
species\_packs:\\
\hspace*{1em}- yaml\_file: data/species\_packs/ar\_2p1\_2p6\_nist.yaml\\
\hspace*{1em}- yaml\_file: data/species\_packs/o\_777\_reduced.yaml\\
\hspace*{2em}namespace: oxygen
\end{quote}
この構成では、同じ領域のArとOが共通の電子密度とEEDFを使い、それぞれの反応係数、供給密度、放射分岐から発光を計算する。Oに対するO$_2$消光のように、混合相手を反応の衝突相手として指定することもできる。

pack内の相対パスはpack自身を基準に解決する。\code{namespace}はpack内で宣言された状態・反応等のIDとその内部参照へ接頭辞を付ける。pack外の供給粒子・衝突相手は共通密度キーを参照するため、異なるpackを同じ組成条件へ接続できる。重複IDは入力時に検出する。新版への置換は旧packの選択を外して新版を指定する手順とし、使用したデータ集合を一意に定める。

\begin{table}[htbp]
\centering\small
\begin{tabularx}{\linewidth}{p{25mm}YY}
\toprule
対象例 & 現在利用できる構成 & 適用条件を整えるための入力 \\
\midrule
Ar & 2p1/2p6の放射分岐pack、六上準位の直接励起pack & 対象条件に応じた準安定密度、cascade・消光・捕獲の採否 \\
O/O$_2$ & O 777 nmの励起とO$_2$消光、Arとの合成例 & O密度と、対象条件に対応する状態別励起断面積 \\
Cl/Cl$_2$ & Cl有効発光、Arとの共有状態例、Cl$_2$帯形 & Cl供給密度、有効励起曲線、822 nmの解離励起寄与 \\
F/NF$_3$ & F線と分子帯を含む多視線の生成例 & F密度、帯の有効供給密度、励起・消光のデータ \\
CF$_4$/O$_2$/Ar & Ar/O packとCF、CF$_2$等の帯の合成例 & 各ラジカル密度と帯の係数・収率・profile \\
BCl$_3$/Ar等 & 追加caseを記述する設定ひな形 & 状態・反応・断面積・放射データの具体化 \\
\bottomrule
\end{tabularx}
\caption{現在のガス別構成と利用時の準備。Arの六上準位には外部断面積候補、O/Cl等の例には例示・有効データを用いる。検証結果は第\ref{sec:benchmarks}節、実測への展開は第\ref{sec:outlook}節に示す。}
\label{tab:gases}
\end{table}

混合ガスの計算では、親ガス分率、ラジカル・準安定密度、EEDFを入力として与える。各種の密度は一つの入力欄で定義し、解く準位密度と固定供給密度の対応を明示する。ガス添加による電子状態や解離組成の変化は、条件ごとにこの入力を更新して表現する。外部の電子・化学モデルを用いる場合も、その結果を領域ごとのEEDFと供給密度へ対応づけられる。したがって、同じ分光計算を使いながら、実測に基づく入力と外部計算に基づく入力を比較できる。

\subsection{繰返し計算と結果の再現性}
compile時には、状態参照、格子、反応構造、固有線形状、視線重み、装置の観測計画を準備する。繰返しfitでは状態量に依存する率と準位密度を更新し、固定構造を再利用する。この分離によって、物理項を振る感度解析や多数の順解析を伴う推定を同じ計算経路で実行できる。

入力の正規化、相対パス補正、schemaと意味検査、観測IDと予測の対応、校正基準の照合が、モデルと測定の接続を支える。断面積のhashと成分別の出力単位を保存し、観測、予測、残差をCSVへ出力するため、第三者は採用したデータと結果を対応づけて確認できる。通常の最適化ではスペクトル計算を反復し、最終解の詳細評価を\code{--audit}で追加する。これにより、推定に必要な反復と、結果の解釈に必要な診断を段階的に実行できる。

\clearpage
\section{機能群の連携、計算フローと数値実装}
\label{sec:implementation}
前節までに述べた物理式と推定法は、入力の準備、状態ごとの反復評価、最終結果の出力という三つの段階で実行される。本節では、それぞれの機能が何を受け取り、次の機能へ何を渡すかを整理する。続いて、積分、行列解法、観測格子、誤差処理を具体的に示し、これらを同じ解析で使うための設計上の工夫を説明する。

\subsection{機能群の役割と受け渡す量}
\begin{table}[htbp]
\centering\small
\begin{tabularx}{\linewidth}{p{24mm}YY}
\toprule
機能群 & 入力と主な処理 & 次へ渡す量・主な実装 \\
\midrule
設定・データ & case、種別データ、装置、観測の参照を展開し、単位と対応関係を検査 & 正規化した設定、状態・反応・放射記録。\code{io}, \code{data} \\
物理計算 & 領域別の電子状態と供給密度から率、準位収支、発光を評価 & 準位密度と体積分光発光。\code{physics}, \code{forward.emissivity} \\
視線・装置 & 発光を経路長で投影し、透過率、校正、LSF、画素平均を適用 & 視線・装置別の予測スペクトル。\code{geometry}, \code{instrument} \\
測定・推定 & 有効画素と波長窓を選び、共分散で残差を白色化して状態量を更新 & 名前付き推定量、予測・測定・残差。\code{inverse} \\
診断・評価 & 最終点の収支と数値品質、感度、profile、保留観測を評価 & 推定を支える情報とモデル改訂の根拠。\code{inverse}, \code{analysis} \\
\bottomrule
\end{tabularx}
\caption{計算機能の分担。\code{forward.compiled}が固定構造を準備し、\code{forward.model}が物理計算から観測予測までを接続する。各名称は\code{oescr}内の相対名である。}
\label{tab:function-groups}
\end{table}

設定・データの段階では、反応の参照先、密度を与える欄、装置IDと観測の対応を確定する。物理計算は、その設定を使って領域ごとの発光を返す。幾何と装置は、同じ発光を実験の視線・波長画素へ変換する。逆解析はこの予測値を測定と比較し、指定した状態量を更新する。この受け渡しを保つことで、例えば消光係数の改訂は準位密度と発光へ、LSFの改訂は測定予測へ、それぞれ対応する段階から反映される。

\subsection{順解析と逆解析の実行フロー}
図\ref{fig:computation-flow}に全体の流れを示す。順解析では、与えた状態に対して発光・観測の経路を一度評価する。逆解析では、その経路を目的関数として繰り返し呼び出し、状態量を測定に合わせる。

最初に、物理・装置の固定構造と、測定の画素状態・波長窓・共分散を準備する。次に、試行するパラメータから領域別のEEDFと密度を構成し、率積分、CR、発光、投影、装置応答を順に計算する。得られた予測と測定の残差を使って、最適化器が次の試行値を決める。終了解では物理量とスペクトルを標準形式で保存し、必要に応じて感度と数値品質を詳しく点検する。

\begin{figure}[p]
\centering\includegraphics[width=\linewidth]{\figfile{computation_flow.pdf}}
\caption{固定構造の準備、反復評価、最終出力を分けた計算フロー。左側は物理・観測の予測経路、右側は測定と推定の準備を表す。逆解析では、残差に基づくパラメータ更新が電子状態・密度の評価へ戻る。順解析は、この予測経路を指定した状態で一度評価する。最終診断は、推定結果を受けて実行する。$k$は反応速度係数、$\boldsymbol r$は白色化残差を表す。}
\label{fig:computation-flow}
\end{figure}

\subsection{固定構造の準備と反復時の再利用}
温度や電子密度の推定では、同じ状態集合、反応集合、視線、分光器を使って多数のスペクトルを計算する。そこで、\code{compile\_case}は入力の展開と参照検査をまとめ、\code{CompiledCase}に固定構造を保持する。状態の索引、反応の組立て規則、エネルギー・波長格子、原子線の基底、分子帯の輪郭、視線重み、観測格子への写像と画素積分の計画をこの段階で準備する。

反復時には、\code{plasma\_state}や\code{eedf}の値を更新し、EEDF、反応率、準位密度、発光振幅を再計算する。断面積を固定エネルギー格子へ補間した値は率計算器が保持する。反応、状態、幾何、装置等の構造を変更した場合は、対応する計算計画を再構築する。この扱いにより、係数やモデルを改訂した結果と、固定量の再利用とを整合させる。

\begin{table}[htbp]
\centering\small
\begin{tabularx}{\linewidth}{p{29mm}YY}
\toprule
処理の時期 & 主に準備・更新する量 & 数値計算上の役割 \\
\midrule
固定構造の準備 & 状態・反応の対応、格子、線・帯の形、視線と観測の写像 & 繰返しに共通する処理をまとめる \\
各試行状態の評価 & EEDF、外部密度、率、準位密度、発光、測定予測 & 状態量に依存する物理応答を更新する \\
最終点の評価 & 詳細収支、品質診断、感度、profile、局所曲率 & 得られた推定値の解釈と次の測定を支える \\
\bottomrule
\end{tabularx}
\caption{処理の頻度に応じた計算の分離。温度・密度の反復評価と、反応集合や装置を変更する再構成の役割を明確にする。}
\end{table}

\subsection{速度係数積分と準位系の数値解法}
エネルギー格子$E_q$上で、式\eqref{eq:rate}の積分核を$g_{r,q}=\sigma_r(E_q)v(E_q)\fe(E_q)$とすると、実装の台形則は
\begin{equation}
k_r^{(h)}=\sum_{q=0}^{N_E-2}\frac{E_{q+1}-E_q}{2}
\left(g_{r,q}+g_{r,q+1}\right)
\label{eq:discrete-rate}
\end{equation}
となる。断面積は固定格子への補間値を再利用し、電子速さとEEDFをその格子上で評価する。分布の更新によって積分核を変える。格子を細分化した比較は離散化への感度を、積分上限を広げた比較は高エネルギー成分への感度を調べる。これらを分けて評価すると、率の変化を格子間隔と積分範囲のどちらに帰属できるかを整理しやすい。

CRでは、各反応・放射・壁損失を一つずつ行列と右辺への寄与として組み立て、総和から式\eqref{eq:matrix}を作る。領域ごとに密行列の直接解法を適用するため、準位数$N_{\mathrm{lev}}$に対する主要な解法費用は概ね$N_{\mathrm{lev}}^3$に従う。選択準位を扱う縮約モデルに適した構成であり、追跡準位数は必要な励起・cascade経路に応じて決める。

同じ反応寄与から準位ごとの生成・損失収支を出力することも、本実装の特徴である。行列の係数、定常密度、反応ごとの収支を相互に照合でき、消光や段階励起を加えたときの作用を説明できる。数値品質の評価には、解の一意性、条件数、相対残差、非負化前の密度を用いる。

\subsection{発光格子、視線一括演算と画素積分}
固有発光では、原子線の解析的な規格化と、分子帯の元の支持域での規格化を準備する。領域ごとの密度と放射率で振幅を更新し、各成分を内部の細格子へ合成する。格子間隔は、固有線幅、帯の表間隔、LSF幅、画素間隔のうち、細かい表現が必要な尺度に合わせて選ぶ。

狭い原子線だけを扱うケースでは、観測装置の波長範囲と余白の和集合を使った分割格子を構成する。離れた観測帯域の間にある空白域の格子量を抑えられる。分子帯を含む場合や広い原子線の場合は、帯・線の支持域を保つ連続格子を使う。したがって、格子の選択は発光成分と観測帯域の両方に基づく。

領域発光を並べた配列へ視線行列$W$を掛けることで、複数視線を一括して投影する。同じ分光器を共有する視線には、波長軸を共通にした透過率・校正・LSFを適用する。LSFは各観測評価で核を作り、波長方向だけにFFT畳み込みを行う。この配列構成により、物理的な視線の違いと共通の装置応答を自然に扱える。

最後に、装置応答後の細格子値$y_q$を区分線形の関数として画素幅で積分する。$h_q=\lambda_{q+1}-\lambda_q$、$0\leq\xi\leq h_q$とすると、累積面積$H$は区間内で
\begin{equation}
H(\lambda_q+\xi)=H(\lambda_q)+\xi y_q+
\frac{\xi^2}{2h_q}\left(y_{q+1}-y_q\right)
\label{eq:antiderivative}
\end{equation}
と評価できる。各画素の境界$e_i,e_{i+1}$でこの累積面積を求め、
\begin{equation}
F_i=\frac{H(e_{i+1})-H(e_i)}{e_{i+1}-e_i}
\label{eq:bin-average-numeric}
\end{equation}
を出力する。画素境界が細格子点の間に来る場合にも、その部分区間の面積を含められる。未分解線や非等間隔画素を扱う際に、画素中心の点標本化に起因する積分誤差を抑え、線中心の移動に伴う各画素の光量変化を評価する処理である。

\subsection{共分散と補正量を共有した目的関数評価}
測定の準備では、有効画素、推定用の波長、観測と予測の対応を固定する。一般の測定共分散には、選択画素に対応するCholesky因子を保持して三角解法を用いる。装置gain、offset、tiltを推定する場合にも同じ画素と白色化を使い、物理量と補正量を一つの誤差基準へそろえる。

共通校正の乗算誤差には、式\eqref{eq:calcov}のrank-one構造を利用する。一つの校正groupについて、独立側の共分散を$C_0=LL^{\mathsf T}$、その白色化を$B=L^{-1}$、校正ベクトルを$\boldsymbol u$とする。$\boldsymbol w=B(\boldsymbol F-\boldsymbol y)$、$\boldsymbol v=B\boldsymbol u$、$a=\sqrt{1+\boldsymbol v^{\mathsf T}\boldsymbol v}$とおけば、補正後残差は
\begin{equation}
\boldsymbol r=\boldsymbol w-
\frac{\boldsymbol v\left(\boldsymbol v^{\mathsf T}\boldsymbol w\right)}{a(a+1)}
\label{eq:rank-one-whitening}
\end{equation}
と書ける。これは$(I+\boldsymbol v\boldsymbol v^{\mathsf T})^{-1/2}\boldsymbol w$の評価に当たる。共通の尺度誤差が繰返し測定へ及ぼす相関を保ちつつ、その追加のために全画素の密な共分散を再分解する費用を抑える。

比や帯面積も使う場合は、選択したfeatureとその共分散を目的関数へ組み込む。スペクトルだけで補正量が線形なら重み付き線形最小二乗を、featureや事前項を含む場合には対応する内側最適化を使う。この内側評価を含めた残差を外側へ返すことで、温度・密度に対する感度も実際に使う目的関数から求められる。

\subsection{最適化と最終診断の分離}
密度等の正の量は$\log_{10}$座標を選べるため、桁の異なる量を有限範囲へ写して推定できる。広い範囲を探索するときはdifferential evolutionで候補点を求め、その後に境界付き最小二乗で精密化する。局所推定だけを使う構成も選べる。計算時間と探索範囲は、推定量数と初期値についての知識に応じて設定する。

反復中は予測スペクトルと残差の計算を中心に実行する。実行時に\code{--audit}を指定すると、終了解で順解析品質、測定だけのJacobian、特異値、条件付きprofile、局所曲率を求める。測定残差数がパラメータ数以上のJacobianにはeconomy SVDを使い、観測空間の大きな正方行列の構成を抑える。パラメータ数が測定残差数を上回る場合は、右零空間の方向も保持する。

この分離は、反復計算の作業量を抑えるだけでなく、結果を読む順序も明確にする。まず測定との一致を確認し、次に観測が支えるパラメータの組合せを調べ、その後に事前分布を含む局所不確かさを評価する。複数の線・視線を追加する効果も、この同じ順序で比較できる。

\subsection{本実装の工夫と検証との対応}
\begin{table}[htbp]
\centering\small
\begin{tabularx}{\linewidth}{p{34mm}YY}
\toprule
設計上の工夫 & 対象となる数値上の課題 & 確認方法 \\
\midrule
固定構造と状態更新の分離 & 繰返しの準備費用、構造変更との整合 & 計画再利用と構造変更時の再構築を確認 \\
発光基準の規格化と細格子 & 観測区間の変更、未分解線の積分 & 解析光量、Gaussian線と画素位置の比較 \\
過程寄与からの行列・収支生成 & 供給・損失の符号と表示の整合 & 二準位・cascadeの解析解、寄与の総和 \\
共通白色化と校正相関の利用 & 物理量とgainの重みの整合、共通尺度誤差 & 直接構成した共分散による評価と比較 \\
推定と測定感度の段階的評価 & 状態量の縮退、事前条件の影響の把握 & 既知の縮退問題、生成回復、profile \\
\bottomrule
\end{tabularx}
\caption{計算の見通しと一貫性を保つための実装上の選択。解析解、保存則、共分散の一致、既知の感度問題を使って、それぞれの性質を確認する。}
\label{tab:numerical-design}
\end{table}

これらの工夫は、物理データを変更した際の影響を反応収支から測定残差まで追跡し、同じ観測基準で推定結果を比較するために組み合わせている。第\ref{sec:benchmarks}節では、この計算経路を具体的な問題設定に適用し、積分量と状態量の回復を評価する。

\clearpage
\section{ベンチマークの問題設定と再計算結果}
\label{sec:benchmarks}
\subsection{検証設計と評価量}
本節は2026年10月27日の作業ツリーから固定したソースと入力を用いた再実行である。未コミット変更を含むため、Gitのcommit番号だけで実行対象を同定せず、同梱の\code{source\_snapshot.zip}と\code{source\_snapshot\_manifest.json}で対象を固定する。生成図の数値、標準result、実行ログ、環境・hashは\code{evidence/}に保存した。原子・分子データが実際の外部資料に基づく場合でも、観測側を同じモデルから生成したケースは独立実験検証には数えない。

パラメータの真値$\theta_0$に対する回復誤差と、スペクトルの相対L2差を
\begin{equation}
\epsilon_\theta=\frac{|\widehat\theta-\theta_0|}{|\theta_0|},\qquad
\epsilon_S=\frac{\|\boldsymbol F-\boldsymbol y\|_2}{\|\boldsymbol y\|_2}
\label{eq:metrics}
\end{equation}
とする。共通状態runnerのfit/holdout指標は、宣言した各チャネルの全21画素を連結して式\eqref{eq:metrics}を求める。補助CSVでは実際のfit maskに選ばれた18画素に限定した値も出す。使用する画素範囲が異なるため同じ数値にはならず、本節の共通状態の表はチャネル全画素によるrunner値で統一した。

標準化RMSEは、選択した画素の$(F_i-y_i)/\sigma_i$の二乗平均平方根である。相対L2、標準化RMSE、事前項等を含む最適化costは異なる量で、互いに代替しない。ベンチマークのPASSは、個別に定めたソフトウェア回復条件を満たす意味である。

\begin{table}[htbp]
\centering\small
\begin{tabularx}{\linewidth}{p{25mm}YYp{28mm}}
\toprule
ケース & 問い & 観測・推定対象 & 根拠の種類 \\
\midrule
解析解・保存 & 次元換算と積分が整合するか & 500 nm帯と600 nm原子線、1領域 & 解析解との比較 \\
二帯域 & 比と絶対尺度を使い分けられるか & 対象密度または電子密度を各1変数 & 雑音なし生成回復 \\
共有状態Ar/O$_2$ & 複数狭帯域が共通状態を制約するか & 5観測、$\Te$と$\nelec$ & 雑音付き生成回復 \\
共有状態Ar/Cl$_2$ & 小残差が回復を保証するか & 7観測、$\Te$と$\nelec$ & 雑音付き生成回復 \\
六kernel Ar & 広帯域入力と別上準位の保留評価が接続するか & 2チャネル各3反復、$\Te$のみ & 同一モデル生成回復 \\
NF$_3$/Ar、Cl$_2$/Ar & 多視線と分子帯の推定が動くか & 5視線、3殻の選択状態量 & 正則化付き生成回復 \\
\bottomrule
\end{tabularx}
\caption{本稿で再実行したケース群。外部実測に対する温度・密度の定量精度を示すケースは、この表にはない。}
\end{table}

\subsection{解析的な原子線・分子帯の積分確認}
仮想粒子X、$p=1\unit{Pa}$、$T_g=300\unit{K}$、$\nelec=3\times10^{16}\unit{m^{-3}}$、半径0.1 mの一様領域を考える。500 nmのGaussian帯はFWHM 1 nm、$k_b=2\times10^{-14}\unit{m^3s^{-1}}$、$Y_b=0.6$、$B_b=0.5$とする。600 nmの原子線は、励起率$10^{-15}\unit{m^3s^{-1}}$、$A=10^7\unit{s^{-1}}$、固有FWHM 0.05 nmで、他の損失を置かない。視線は中心を通り、光路長0.2 mである。装置のFWHMは0.05 nm、画素幅は0.05 nm、透過率1、恒等の絶対校正を使う。

原子線の積分発光は式\eqref{eq:single}と\eqref{eq:line}から厳密に計算できる。分子帯は中心波長の光子エネルギーを使う狭帯域近似$hc/\lambda_0$を基準とする。コードは$hc/\lambda$を帯内で積分するので、この比較には格子誤差以外に小さい有限幅補正が含まれる。

\begin{table}[htbp]
\centering\small
\begin{tabular}{lrrr}
\toprule
量 & 解析値 & 計算値 & 相対差 \\
\midrule
原子線積分発光 & 190.8237101 & 190.8237101 & $1.11\times10^{-16}$ \\
分子帯積分発光 & 1373.9307124 & 1373.9317035 & $7.21\times10^{-7}$ \\
合計積分発光 & 1564.7544224 & 1564.7554135 & $6.33\times10^{-7}$ \\
中心視線積分輝度 & 312.9508845 & 312.9510827 & $6.33\times10^{-7}$ \\
\bottomrule
\end{tabular}
\caption{発光はW m$^{-3}$ sr$^{-1}$、視線輝度はW m$^{-2}$ sr$^{-1}$。図\ref{fig:analytic}は装置出力と相対差を示す。基準許容値$2\times10^{-4}$より十分小さい。}
\end{table}

\begin{figure}[htbp]
\centering\includegraphics[width=\linewidth]{\figfile{benchmark_analytic.pdf}}
\caption{解析的保存例。(a) 500 nm帯と600 nm線の中心視線スペクトル。(b) 積分量の絶対相対差を百分率で表示。破線は0.02\%の許容値。帯幅が異なるためpeak高さの比較では積分光量は判断できない。数値は表と\code{analytic\_conservation.csv}に対応する。}
\label{fig:analytic}
\end{figure}

\subsection{二帯域の最小例：比、actinometry、絶対尺度}
500 nmの対象帯と510 nmの参照帯を用いる。両者のFWHMは0.6 nm、装置FWHMは0.3 nm、反応係数はそれぞれ$10^{-14}$と$2\times10^{-14}\unit{m^3s^{-1}}$、光子収率と分岐は1である。真値は$N_{\mathrm{Target}}=2\times10^{18}$、$N_{\mathrm{Actinometer}}=10^{18}$、$\nelec=3\times10^{16}\unit{m^{-3}}$とし、$\Te=3\unit{eV}$を固定する。一定係数のため、この例は温度感度を検証しない。

比診断とactinometryでは、同じ面積比を使って対象密度だけを推定した。両モードで$1.9999999997\times10^{18}\unit{m^{-3}}$、相対誤差$1.61\times10^{-10}$を得た。絶対モードでは供給密度を固定し、$\nelec=3.0000000000001\times10^{16}\unit{m^{-3}}$、相対誤差$3.31\times10^{-14}$を得た。いずれも雑音なしの生成入力と同じ式を使うため、極小の誤差は自己整合性試験で期待される結果である。確認しているのは、情報源の切替、単位、設定から結果までの接続である。現実のactinometryの系統誤差をこの数値で評価することはできない。

\subsection{\texorpdfstring{Ar/O$_2$とAr/Cl$_2$}{Ar/O2とAr/Cl2}の共有状態}
両者とも半径0.05 m、一殻、中心視線、$T_g=400\unit{K}$、$E=0$--60 eVを601点で積分する。装置は各線を含む4 nm幅、0.2 nm画素、0.35 nm FWHMで、固定の絶対尺度と1\%の校正標準不確かさを置く。生成測定には1.5\%の点雑音とpeakの0.15\%を下限尺度とする雑音を加えた固定fixtureを使う。両ケースとも$\Te$と$\log_{10}\nelec$を推定し、放射捕獲と壁損失を置かない。

Ar/O$_2$は$p=2\unit{Pa}$、親ガス分率0.8/0.2、$N_{\mathrm O}=2\times10^{19}\unit{m^{-3}}$を与え、750.4、763.5、777.4 nmをfit、800.6と922.4 nmを保留する。fitにはAr二上準位とO一上準位の三kernelがある。二つの保留Ar線は既にfitした2p6からの分岐で、別の励起kernelではない。

Ar/Cl$_2$は$p=2.4\unit{Pa}$、分率0.78/0.22、$N_{\mathrm{Cl}}=2\times10^{19}\unit{m^{-3}}$を与える。750.4、763.5、800.6、922.4 nmのAr四線をfitするが、独立な上準位は二つである。725.7、754.7、822.2 nmのCl系は例示的な有効経路として保留し、温度・密度fitには使わない。

\begin{table}[htbp]
\centering\small
\begin{tabular}{llrrr}
\toprule
ケース & パラメータ & 真値 & 初期値 & 推定値（相対誤差） \\
\midrule
Ar/O$_2$ & $\Te$ / eV & 3.5 & 2.8 & 3.50596（0.170\%） \\
 & $\nelec$ / $10^{16}$ m$^{-3}$ & 2.0 & 3.0 & 1.96832（1.584\%） \\
Ar/Cl$_2$ & $\Te$ / eV & 4.0 & 3.2 & 3.74073（6.482\%） \\
 & $\nelec$ / $10^{16}$ m$^{-3}$ & 1.6 & 2.4 & 2.07902（29.939\%） \\
\bottomrule
\end{tabular}
\caption{二つの共有状態ベンチマークの回復結果。数値は単一の固定雑音実現に対するもので、誤差分布や信頼区間ではない。}
\label{tab:commonrecovery}
\end{table}

\begin{table}[htbp]
\centering\small
\begin{tabular}{lrr}
\toprule
指標 & Ar/O$_2$ & Ar/Cl$_2$ \\
\midrule
fit測定に対する相対L2（チャネル全画素） & 1.274\% & 0.944\% \\
保留窓の雑音なし真値に対する相対L2 & 0.860\% & 12.297\% \\
保留測定に対する相対L2 & 2.410\% & 12.880\% \\
測定のみの局所rank / 対象数 & 2/2 & 2/2 \\
測定Jacobian条件数 & 15.565 & 220.195 \\
エネルギー格子細分化の相対L2差 & 0.0697\% & 0.0582\% \\
波長格子細分化の相対L2差 & 0.00560\% & 0.00562\% \\
既定の回復判定 & PASS & FAIL \\
\bottomrule
\end{tabular}
\caption{収束、回復、保留予測を別に示した評価。既定ゲートはoptimizer成功、局所full rank、温度・密度誤差と保留真値差がそれぞれ5\%未満、格子細分化差1\%未満を要求する。}
\label{tab:commonmetrics}
\end{table}

\begin{figure}[p]
\centering\includegraphics[width=\linewidth]{\figfile{benchmark_common_state.pdf}}
\caption{共有状態の結果。(a,c) Ar/O$_2$、Ar/Cl$_2$の初期値と推定値を真値で割って表示。点の注記は絶対相対誤差。(b,d) 測定のみの条件付き温度profileで、他の推定量を再最適化する。縦軸は$\Delta\sum r_i^2$で、ゼロ近傍が線形、大きい値が対数の表示、縦破線は真値。(e,f) fitしたO線と保留したCl線を各測定peakで規格化して比較。error barは画素の$\sigma$のみで共通校正誤差を含まない。粗いprofile点の結線は谷の詳細な形や信頼区間を与えない。}
\label{fig:common}
\end{figure}

Ar/Cl$_2$はoptimizerが成功し、局所rankも2だが、回復ゲートを満たさない。Arの二上準位の応答で温度と密度を分ける感度が弱く、温度のずれを密度の増加で補えるため、小さいfit残差と大きい密度誤差が併存する。図\ref{fig:common}のprofileと条件数はその状況と整合する。ただし、この一回の結果だけから弱識別性と雑音以外の寄与を厳密に分解したとはいえない。Clの有効経路には独立した適用性評価もないため、保留の不一致を実験中のCl動力学の失敗率と解釈してはならない。

\subsection{六kernel Ar：広帯域と別上準位の予測確認}
純Ar、$p=2\unit{Pa}$、$T_g=400\unit{K}$、一殻、中心視線、$\nelec=2\times10^{16}\unit{m^{-3}}$固定、真値$\Te=4.0\unit{eV}$、初期値2.5 eVとする。650--850 nmと840--1050 nmの二チャネルに0.2 nm画素を置き、各三回の決定論的摂動を加えた観測を用意する。六つの異なる上準位2p10、2p9、2p8、2p7、2p6、2p1の基底状態からの直接励起断面積は外部公開bundle\cite{schmalzried}、17の放射分岐はNIST ASD\cite{nist}で固定する。16分岐が装置範囲に入り、1148.8 nm分岐も占有率の損失には残す。

五つの上準位をfitに使い、2p10を別kernelとして保留する。fitに使用した上準位の別分岐も保留するが、これは二種類の確認である。相対形状モードで装置ごとのgainを許し、推定する物理量は共通$\Te$一つだけとする。六つの観測CSVが六つの独立状態を意味するわけではない。

推定値は3.983264 eV、真値に対する相対誤差0.4184\%となった。標準化RMSEは全保留窓0.5124、別kernel保留0.4458、分岐保留0.5500で、ソフトウェア回復条件を満たした。真値4 eVでの$d\log k/d\log\Te$は2.854--4.009であり、六つの率が同じ温度関数の定数倍だけではないことを確認できる。ただし、これだけで実験上の独立情報量を確定できない。

\begin{figure}[p]
\centering\includegraphics[width=\linewidth]{\figfile{benchmark_ar_six_kernel.pdf}}
\caption{六kernel Arの広帯域例。二チャネルの測定・初期予測・最適予測、測定のみの温度profile、および各励起率の対数温度感度を示す。profileは宣言したMaxwell分布族内の一変数診断である。推定温度は自由な分布再構成の結果ではなく、固定電子密度の条件下で推定したMaxwellian有効温度である。}
\label{fig:six}
\end{figure}

このケースでは準安定励起、cascade、消光、捕獲の閉包は行っていない。公開断面積と放射データを使っていても、すべてのkernelの根拠区分は\code{validation\_candidate}である。実測Teの適用性評価は保留される。Arの状態別励起には、直接励起とcascadeを区別した測定\cite{chilton}や、電子動力学とCRを組み合わせた実験比較\cite{arellano}が重要であり、それらの結果を本コードの実証結果へ転用することはできない。

\subsection{\texorpdfstring{NF$_3$/Ar、Cl$_2$/Ar}{NF3/Ar、Cl2/Ar}：多視線と経験的帯の例}
両例は半径0.15 mの三殻を五視線から観測する。NF$_3$/Arは$p=133.322\unit{Pa}$、300--820 nmを0.5 nm画素で、Cl$_2$/Arは$p=2.3998\unit{Pa}$、270--885 nmを0.2 nm画素で観測する。文献\cite{an2023,fuller}は運転条件と線・帯の配置の手掛かりだが、五視線の測定はコードから生成した入力である。実験の視線データを取得して再現したものではない。

NF$_3$/Arでは三殻それぞれの$\Te$、F密度、N$_2$帯用の有効供給密度を推定する。Cl$_2$/ArではCl密度とCl$_2$帯用の有効供給密度を三殻で推定し、温度は固定する。いずれも相対形状、gain/offset、事前分布、殻平滑化を伴い、経験的発光成分を含む。\code{N2\_emit}、\code{Cl2\_emit}は帯を表す供給キーで、全N$_2$、全Cl$_2$密度の独立測定量ではない。

\begin{table}[htbp]
\centering\small
\begin{tabularx}{\linewidth}{p{25mm}p{42mm}p{30mm}Y}
\toprule
ケース & 推定対象 & 殻ごとの回復誤差範囲 & 測定のみの局所診断 \\
\midrule
NF$_3$/Ar & $\Te$（3変数） & 5.04--5.95\% & rank 9/9 \\
 & F密度（3変数） & 8.61--12.76\% & 条件数$8.31\times10^5$ \\
 & N$_2$有効供給（3変数） & 12.52--24.92\% &  \\
Cl$_2$/Ar & Cl密度（3変数） & 10.61--20.08\% & rank 6/6 \\
 & Cl$_2$有効供給（3変数） & 6.23--17.14\% & 条件数32.71 \\
\bottomrule
\end{tabularx}
\caption{多視線例の回復結果。事前分布や平滑化は最適化には含むが、ここに示すrankと条件数には含まない。完全rankであっても、NF$_3$例の条件数は大きい。}
\label{tab:multichord}
\end{table}

\begin{figure}[p]
\centering\includegraphics[width=\linewidth]{\figfile{benchmark_multichord.pdf}}
\caption{多視線例の代表スペクトルと半径方向の回復。上段はNF$_3$/Ar、下段はCl$_2$/Ar。左列は中心視線の広帯域比較、右列は三殻の推定量を真値で割った値。右列縦軸は差を読める範囲に限定しておりゼロ始まりではない。スペクトルの一致と個々の供給密度の回復を別に評価する。経験帯、相対gain、事前分布を含むため、この図を実測の絶対密度精度の証明に使うことはできない。}
\label{fig:multichord}
\end{figure}

保存された補正後予測を全波長で連結した相対L2差はNF$_3$/Arで14.52\%、Cl$_2$/Arで18.32\%だった。選択したスペクトル画素だけでは10.78\%、11.63\%である。これは窓featureや事前項を含むoptimizer costとは別の表示用指標である。これらの例は、複数成分、幾何、装置、正則化付き推定を接続する機能を確認するが、各ラジカル密度の高精度回復を示したケースではない。

\subsection{現在の証拠から言えること}
解析例は、発光単位、放射分岐、投影と積分の一貫性を支持する。二帯域と六kernel例は、宣言した情報源からの条件付き回復を支持する。共通状態例は、kernel選択によって温度・密度の分離が大きく変わることを示す。多視線例は複合入力から推定までの処理を確認する。生成と推定に同一の演算子を用いるため、共通する原子データの誤りや欠落機構の影響を、この試験だけで検出することはできない。外部診断に対する定量検証は、次節の課題として残す。

\section{考察と利用範囲}
\subsection{コードの有用性が現れる場面}
第一に、発光の機構と観測の機構を同じ計算で追跡できる。ある線の不足を見たとき、励起率、準安定密度、消光、分岐、光路、透過率のどこを変更したかが明確になる。生成・損失収支と成分別出力は、経験的な係数調整だけで説明を終えず、仮定の影響を調べる手段になる。

第二に、物理項をデータで改訂できる。断面積の新版、未観測の放射分岐、追加の消光相手を、species packや反応集合の変更として扱える。原子データを直す作業と、励起の閉包を変更する作業を分けられるため、複数ガスで同じAr参照を再利用しやすい。関数形が変わればplugin、線形の前提が変わればsolver変更、と作業の範囲も判断しやすい。

第三に、観測設計を定量的に検討できる。分岐線を増やすこと、別上準位を追加すること、別種を追加すること、視線を追加することは、それぞれ増える情報が違う。Ar/Cl$_2$の例では、四本のAr線でも上準位が二つなら温度・密度分離が弱い場合がある。測定Jacobianとprofileを使うことで、誤差を減らすべき装置条件と、独立に測るべきプラズマ量を検討できる。

これらは、適用範囲を明示した分光診断の試作、既存CRMの縮約版の比較、校正と視線条件の感度調査、共同研究でのデータ・反応集合の交換に適している。放電の電力・流量から最終的な組成や加工速度を直接予測する課題では、外部の化学・輸送・表面モデルが別に必要になる。

\subsection{外部実測による検証へ進む条件}
現在の外部候補には、Ar-CCPの線比資料\cite{arellano}と、低圧N$_2$/O$_2$のNO紫外発光資料\cite{schuecke}がある。前者には電子分布、準安定・cascade、断面積の整合性、後者には主要なN$_2$(A)密度など、必要入力の閉包が残る。参考論文中のCRMやPIC/MCCが実験と一致したことは、別の反応集合を持つOESCRの検証ではない。

実測へ進める単位は、gas名一般ではなく、反応集合・遷移・圧力範囲・観測幾何を指定した一つのケースである。校正済みスペクトル、波長応答、LSF、光路、ガス温度、供給粒子密度とその不確かさをまとめ、温度または電子密度の独立参照を対応づける。fitに使わない波長・条件で予測を確認し、合否値は結果を見る前に定める。

BOLSIG+等の固定出力とのEEDF・速度係数比較は、分布の読込と積分の外部数値確認に有用である\cite{hagelaar,lxcat}。ただし、現在のtemplateと実行準備の存在を、独立実行結果との一致と取り違えない。その一致が得られても、光学校正や準位閉包を含む実験的Te精度を単独では保証しない。

Fullerらのactinometry\cite{fuller}やNF$_3$洗浄の分光研究\cite{an2023}は、対象種と実験条件の設計に有益である。一方、同じOESから導出された電子温度やactinometry密度を、そのOESモデルの独立真値として再利用すると循環評価になる。図のdigitization誤差、実験誤差、モデル不一致は区別して扱う必要がある。

\subsection{主要な近似と改善の優先順位}
物理的な限界は、定常・選択準位・線形衝突収支、外部密度の固定、領域間輸送の省略、有効分子帯、入力光学厚さによる簡略捕獲、有効壁損失に集約できる。逆問題の限界は、低次元EEDF族への条件づけ、局所SVDと局所profile、曲率に基づく不確かさである。これらは異なる原因なので、単に自由パラメータを増やして同時に解消できるわけではない。

改善では、まず一つの校正済み実測ケースについて供給密度と励起経路を閉じ、独立診断と保留スペクトルで誤差を測る。その結果に基づいて、準安定励起、cascade、消光、捕獲の必要性を判断する。一般のEEDF再構成や複雑な化学連成を先に増やすと、観測が支えない自由度を増やす可能性がある。現行の監査機能は、こうした拡張の前に何が不足しているかを見定めるためにも有用である。

\subsection{結論}
OESCRは、明示した電子状態と粒子密度から、選択準位の線形CR、発光、視線幾何、装置応答、測定残差までを一貫して計算する。物理項の所在と単位が分かれ、複数種のデータを合成できる点は、モデルの点検と改訂を容易にする。再現例は、この計算経路と条件付き逆推定の機能を支持する。同時にAr/Cl$_2$の回復失敗と多視線例の残る誤差は、スペクトル適合だけで物理量を確定できないことを示した。第三者が利用する際は、指定した反応・入力・観測範囲に対する診断モデルとして扱い、実験的な定量主張を独立検証のある範囲へ限定することが妥当である。

\clearpage
\appendix
\section{主要記号と実装との対応}
\small
\begin{longtable}{p{24mm}p{43mm}p{86mm}}
\toprule 記号 & 意味・単位 & 主な設定または実装 \\
\midrule\endhead
$\Te,E_D,T_c,T_h$ & EEDFの温度・尺度 / eV & \code{eedf.kind}, \code{eedf.zones}、\path{physics/eedf.py} \\
$\alpha$ & 高温電子の粒子割合 & \code{hot\_fraction} \\
$\fe$ & エネルギー確率密度 / eV$^{-1}$ & \code{build\_pdf}。有限格子で正規化 \\
$\nelec$ & 電子密度 / m$^{-3}$ & \code{plasma\_state.ne\_shells\_m3} \\
$p,T_g,x_s$ & 圧力 / Pa、温度 / K、分率 & \code{gas\_mixture}、\path{physics/residuals.py} \\
$N_s$ & 外部供給粒子密度 / m$^{-3}$ & \code{radicals}, \code{metastables}, \code{residuals}等 \\
$\sigma_r,k_r$ & 断面積 / m$^2$、率 / m$^3$ s$^{-1}$ & \code{cross\_section\_file}、\path{physics/rates.py} \\
$n_i^*,M,\boldsymbol b$ & 準位密度、損失行列、供給 & \code{states}, \code{reactions}、\path{physics/cr_atomic.py}, \path{physics/cr_processes.py} \\
$A,\beta,\tau$ & 放射率、脱出係数、光学厚さ & \code{transitions}, \code{trapping} \\
$\gamma,f_{\rm flux},L$ & 壁確率、流束係数、長さ / m & \code{wall}、\path{physics/wall.py} \\
$j_\lambda$ & 体積分光発光 & \path{forward/emissivity.py}, \path{physics/bands.py} \\
$W,I_\lambda$ & 経路長 / m、分光放射輝度 & \code{geometry}、\path{geometry/axisym_shell.py} \\
$K,T,\mathcal C$ & LSF、透過率、単位換算 & \code{instruments}、\path{instrument/spec.py}等 \\
$\boldsymbol F,\boldsymbol y$ & 予測と測定画素平均 & \code{prediction\_key}, 観測CSV \\
$C,J,D_p$ & 共分散、測定感度、条件付きprofile & \path{inverse/measurements.py}, \path{inverse/identifiability.py}, \path{inverse/profiles.py} \\
\bottomrule
\end{longtable}
\normalsize

状態の\code{energy\_eV}と\code{degeneracy}は、準位の同定と原子データの記録に用いる。速度係数を決める閾値と断面積は反応レコード、光子エネルギーを決める波長は遷移レコードで管理する。したがって、準位データの改訂時には、対応する反応・遷移の値を併せて確認する。詳細釣合いに基づく脱励起率を採用する場合は、その関係を率モデルとして明示する。

\section{再現手順と添付データ}
\small
本文の数値は付属の\code{evidence/}にあるCSV・YAML・JSONと対応する。通常の作業ツリーに対しては、リポジトリルートから次の入口を使える。
\begin{quote}\small\ttfamily
python scripts/run\_common\_state\_benchmarks.py --out runs/common\\
python scripts/run\_argon\_six\_kernel\_candidate.py --out runs/ar6\\
python scripts/run\_project.py project.yaml --out runs/fit --audit
\end{quote}
二帯域は\path{examples/use_cases/two_band/}の初期caseと三つのinverse設定、多視線は\path{examples/benchmarks/}の各\code{project.yaml}を使う。common-state runnerの終了値1は、今回のAr/Cl$_2$が所定の回復基準を超えたことを表す。この結果も、推定誤差と感度を考察する対象として本文に掲載した。

厳密に掲載時点を再現する場合は、\code{source\_snapshot.zip}を専用の空ディレクトリへ展開し、そのルートを作業ディレクトリとして依存を導入する。archiveのSHA-256先頭16桁は\code{16a5f0c95cf381c2}で、完全な値はmanifestにある。付属再現スクリプトと\code{run\_manifest.json}が、完全な実行引数、種、hash、結果位置を保持する。再現時には、同じsnapshot内のコードと入力を一組として用いる。図は\code{figures/}のベクトルPDF、数値は\code{evidence/}のCSVで確認できる。

再計算環境はWindows 11、Python 3.13.12、NumPy 2.4.3、SciPy 1.17.1、Matplotlib 3.11.2である。
\normalsize

\clearpage
\phantomsection
\addcontentsline{toc}{section}{参考文献}
\begingroup
\small\raggedright
\begin{thebibliography}{99}

\bibitem{fujimoto}
T.~Fujimoto,
\textit{Plasma Spectroscopy},
Oxford University Press (2004).
\href{https://doi.org/10.1093/acprof:oso/9780198530282.001.0001}{doi:10.1093/acprof:oso/9780198530282.001.0001}.

\bibitem{lieberman}
M.~A.~Lieberman and A.~J.~Lichtenberg,
\textit{Principles of Plasma Discharges and Materials Processing},
2nd ed., Wiley (2005).
\href{https://doi.org/10.1002/0471724254}{doi:10.1002/0471724254}.

\bibitem{coburn}
J.~W.~Coburn and M.~Chen,
``Optical emission spectroscopy of reactive plasmas: A method for correlating emission intensities to reactive particle density,''
\textit{J. Appl. Phys.} \textbf{51}, 3134--3136 (1980).
\href{https://doi.org/10.1063/1.328060}{doi:10.1063/1.328060}.

\bibitem{hagelaar}
G.~J.~M.~Hagelaar and L.~C.~Pitchford,
``Solving the Boltzmann equation to obtain electron transport coefficients and rate coefficients for fluid models,''
\textit{Plasma Sources Sci. Technol.} \textbf{14}, 722--733 (2005).
\href{https://doi.org/10.1088/0963-0252/14/4/011}{doi:10.1088/0963-0252/14/4/011}.

\bibitem{lxcat}
S.~Pancheshnyi, S.~Biagi, M.~C.~Bordage, G.~J.~M.~Hagelaar,
W.~L.~Morgan, A.~V.~Phelps, and L.~C.~Pitchford,
``The LXCat project: Electron scattering cross sections and swarm parameters for low temperature plasma modeling,''
\textit{Chem. Phys.} \textbf{398}, 148--153 (2012).
\href{https://doi.org/10.1016/j.chemphys.2011.04.020}{doi:10.1016/j.chemphys.2011.04.020}.

\bibitem{nist}
A.~Kramida, Yu.~Ralchenko, J.~Reader, and NIST ASD Team,
\textit{NIST Atomic Spectra Database}, version 5.12 (2024),
National Institute of Standards and Technology, Gaithersburg, MD.
\href{https://physics.nist.gov/asd}{https://physics.nist.gov/asd}
（参照日：2026年10月27日）.
\href{https://doi.org/10.18434/T4W30F}{doi:10.18434/T4W30F}.

\bibitem{chilton}
J.~E.~Chilton, J.~B.~Boffard, R.~S.~Schappe, and C.~C.~Lin,
``Measurement of electron-impact excitation into the $3p^{5}4p$ levels of argon using Fourier-transform spectroscopy,''
\textit{Phys. Rev. A} \textbf{57}, 267--277 (1998).
\href{https://doi.org/10.1103/PhysRevA.57.267}{doi:10.1103/PhysRevA.57.267}.

\bibitem{arellano}
F.~J.~Arellano, M.~Gyulai, Z.~Donkó, P.~Hartmann,
T.~V.~Tsankov, U.~Czarnetzki, and S.~Hamaguchi,
``First-principles simulation of optical emission spectra for low-pressure argon plasmas and its experimental validation,''
\textit{Plasma Sources Sci. Technol.} \textbf{32}, 125007 (2023).
\href{https://doi.org/10.1088/1361-6595/ad0ede}{doi:10.1088/1361-6595/ad0ede}.

\bibitem{schuecke}
L.~Schücke, A.~N.~Torres Segura, I.~Korolov, N.~Bibinov,
P.~Awakowicz, and A.~R.~Gibson,
``NO densities and UV emission over the E--H mode transition in a low-pressure inductively coupled plasma device,''
\textit{Plasma Sources Sci. Technol.} \textbf{34}, 045015 (2025).
\href{https://doi.org/10.1088/1361-6595/adcbd3}{doi:10.1088/1361-6595/adcbd3}.

\bibitem{fuller}
N.~C.~M.~Fuller, I.~P.~Herman, and V.~M.~Donnelly,
``Optical actinometry of Cl$_2$, Cl, Cl$^+$, and Ar$^+$ densities in inductively coupled Cl$_2$--Ar plasmas,''
\textit{J. Appl. Phys.} \textbf{90}, 3182--3191 (2001).
\href{https://doi.org/10.1063/1.1391222}{doi:10.1063/1.1391222}.

\bibitem{an2023}
S.~An and S.~J.~Hong,
``Spectroscopic Analysis of NF$_3$ Plasmas with Oxygen Additive for PECVD Chamber Cleaning,''
\textit{Coatings} \textbf{13}(1), 91 (2023).
\href{https://doi.org/10.3390/coatings13010091}{doi:10.3390/coatings13010091}.

\bibitem{schmalzried}
A.~Schmalzried,
\textit{Bolsig+ format cross section of e-N2, O2, NO, Ar, O, N},
version v1, Zenodo (2023), dataset.
\href{https://doi.org/10.5281/zenodo.8192503}{doi:10.5281/zenodo.8192503}.

\end{thebibliography}
\endgroup

\end{document}
